%=======================================================
\chapter{Texto em várias colunas}

\textsc{Qual o pacote:} \textbf{multicol} (Acesso: \url{https://www.ctan.org/pkg/multicol})

\textsc{O que faz?}: Divide o texto em colunas (Máx. de 10 colunas.)
\ \\

\begin{codex}{Código multicolunas com valor 3 e Lorem Ipsum}
\begin{multicols}{3}
\lipsum[2]
\end{multicols}
\end{codex}
\ \\

\textsc{Resultado}:
\ \\

\begin{multicols}{3}
	\lipsum[2]
\end{multicols}

\begin{codex}{Código multicolunas com valor 5 e Lorem Ipsum}
\begin{multicols}{5}
\lipsum[2]
\end{multicols}
\end{codex}
\ \\

\textsc{Resultado}:

\begin{multicols}{5}
	\lipsum[2]
\end{multicols}
\ \\
Também é possível incluir uma linha separadora entre as colunas. Ver a seguir.
\ \\

\begin{codex}{Código multicolunas com valor 3 e Lorem Ipsum + Linha}
\begin{multicols}{3}
\setlength{\columnseprule}{0.2pt}
\lipsum[2]
\end{multicols}
\end{codex}
\ \\

\textsc{Resultado}:

\begin{multicols}{3}
	\setlength{\columnseprule}{0.2pt}
	\lipsum[2]
\end{multicols}

Veja também \url{https://www.alessandroduarte.com.br/?page_id=602} para um tutorial em português.